% GENERATED FILE. Edit canonical lessons, then run npm run generate:books.
% canonical-source-hash: fnv1a64:cdf7c6ae4069d68b
% canonical-lessons: SA-C423-pravesah, SA-C423-nirgamah, SA-C423-viparitah, SA-C423-antaram, SA-C423-dvarapalah

\chapter{Ways In and Out}
\label{ch:sa-ways-in-and-out}
\hlchaptermodality{423}

\begin{chapteropening}
\textbf{By the end of this chapter:} \emph{I can find the way in, the way out and the gap between.}

Complete the last lesson of chapter 423: I can find the way in, the way out and the gap between.
\end{chapteropening}

\section[\texorpdfstring{\sk{प्रवेशः} (praveśaḥ)}{प्रवेशः (praveśaḥ)}]{\sk{प्रवेशः} (praveśaḥ) — an entrance; entry}

\label{lesson:SA-C423-pravesah}



\begin{quote}
Before the new one: say the Sanskrit for a tunnel, then the Sanskrit for a ditch.
\end{quote}

\subsection*{You'll want to know: \sk{प्रवेशः}}
\textbf{\sk{प्रवेशः}} — \emph{praveśaḥ} — \textquotedblleft{}an entrance; entry\textquotedblright{}.

\begin{quote}
\sk{प्रवेशः} \sk{वामतः} \sk{अस्ति।}
\end{quote}

\begin{quote}
\emph{praveśaḥ vāmataḥ asti.}
\end{quote}

\begin{quote}
The entrance is on the left.
\end{quote}

A masculine noun.

\begin{cousinweb}
\textbf{\sk{प्रवेशः}} (\emph{praveśaḥ}) is built from \emph{pra}, \textquotedblleft{}forward\textquotedblright{}, and the root √viś, \textquotedblleft{}to enter\textquotedblright{}, which also gives \emph{praviśati}, \textquotedblleft{}enters\textquotedblright{}; MW gives \emph{praveśa} as entering, entrance (MW).
\end{cousinweb}

\begin{grammarlens}[title={What you've built}]
One more word for getting around.
\end{grammarlens}

\subsection*{Your turn}
\begin{itemize}
\raggedright
  \item \emph{Say it:} \emph{praveśaḥ}
  \item \emph{Say it:} \emph{praveśaḥ}, once more
  \item \emph{Say it:} \emph{praveśaḥ vāmataḥ asti.} — The entrance is on the left.
  \item \emph{Recall:} say the Sanskrit for a tunnel, then the Sanskrit for a ditch, then say \emph{praveśaḥ} again
\end{itemize}

\begin{tcolorbox}[breakable,colback=teal!4,colframe=teal!35!black,title={Before you move on}]
What does \sk{प्रवेशः} mean? (\textquotedblleft{}An entrance; entry\textquotedblright{}.) Say it once more.
\end{tcolorbox}
\section[\texorpdfstring{\sk{निर्गमः} (nirgamaḥ)}{निर्गमः (nirgamaḥ)}]{\sk{निर्गमः} (nirgamaḥ) — a way out; an exit}

\label{lesson:SA-C423-nirgamah}



\begin{quote}
Before the new one: say the Sanskrit for a ditch, then the Sanskrit for an entrance.
\end{quote}

\subsection*{You'll want to know: \sk{निर्गमः}}
\textbf{\sk{निर्गमः}} — \emph{nirgamaḥ} — \textquotedblleft{}a way out; an exit\textquotedblright{}.

\begin{quote}
\sk{निर्गमः} \sk{दक्षिणतः} \sk{अस्ति।}
\end{quote}

\begin{quote}
\emph{nirgamaḥ dakṣiṇataḥ asti.}
\end{quote}

\begin{quote}
The exit is on the right.
\end{quote}

A masculine noun.

\begin{cousinweb}
\textbf{\sk{निर्गमः}} (\emph{nirgamaḥ}) is built from \emph{nir}, \textquotedblleft{}out\textquotedblright{}, and the root √gam, \textquotedblleft{}to go\textquotedblright{}, as \emph{nirgacchati}, \textquotedblleft{}goes out\textquotedblright{}, is; MW gives \emph{nirgama} as going forth, departure, and the way out of a building is that sense in modern use (MW).
\end{cousinweb}

\begin{grammarlens}[title={What you've built}]
One more word for getting around.
\end{grammarlens}

\subsection*{Your turn}
\begin{itemize}
\raggedright
  \item \emph{Say it:} \emph{nirgamaḥ}
  \item \emph{Say it:} \emph{nirgamaḥ}, once more
  \item \emph{Say it:} \emph{nirgamaḥ dakṣiṇataḥ asti.} — The exit is on the right.
  \item \emph{Recall:} say the Sanskrit for a ditch, then the Sanskrit for an entrance, then say \emph{nirgamaḥ} again
\end{itemize}

\begin{tcolorbox}[breakable,colback=teal!4,colframe=teal!35!black,title={Before you move on}]
What does \sk{निर्गमः} mean? (\textquotedblleft{}A way out; an exit\textquotedblright{}.) Say it once more.
\end{tcolorbox}
\section[\texorpdfstring{\sk{विपरीतः} (viparītaḥ)}{विपरीतः (viparītaḥ)}]{\sk{विपरीतः} (viparītaḥ) — opposite; reversed}

\label{lesson:SA-C423-viparitah}



\begin{quote}
Before the new one: say the Sanskrit for an entrance, then the Sanskrit for an exit.
\end{quote}

\subsection*{You'll want to know: \sk{विपरीतः}}
\textbf{\sk{विपरीतः}} — \emph{viparītaḥ} — \textquotedblleft{}opposite; reversed\textquotedblright{}.

\begin{quote}
\sk{मम} \sk{गृहं} \sk{विपरीतायां} \sk{दिशायाम्} \sk{अस्ति।}
\end{quote}

\begin{quote}
\emph{mama gṛhaṃ viparītāyāṃ diśāyām asti.}
\end{quote}

\begin{quote}
My house is in the opposite direction.
\end{quote}

\begin{cousinweb}
\textbf{\sk{विपरीतः}} (\emph{viparītaḥ}) is built from \emph{vi}, \emph{pari} and the root √i, \textquotedblleft{}to go\textquotedblright{}: gone round the other way; MW gives \emph{viparīta} as turned round, reversed, opposite (MW).
\end{cousinweb}

\begin{grammarlens}[title={What you've built}]
One more word for getting around.
\end{grammarlens}

\subsection*{Your turn}
\begin{itemize}
\raggedright
  \item \emph{Say it:} \emph{viparītaḥ}
  \item \emph{Say it:} \emph{viparītaḥ}, once more
  \item \emph{Say it:} \emph{mama gṛhaṃ viparītāyāṃ diśāyām asti.} — My house is in the opposite direction.
  \item \emph{Recall:} say the Sanskrit for an entrance, then the Sanskrit for an exit, then say \emph{viparītaḥ} again
\end{itemize}

\begin{tcolorbox}[breakable,colback=teal!4,colframe=teal!35!black,title={Before you move on}]
What does \sk{विपरीतः} mean? (\textquotedblleft{}Opposite; reversed\textquotedblright{}.) Say it once more.
\end{tcolorbox}
\section[\texorpdfstring{\sk{अन्तरम्} (antaram)}{अन्तरम् (antaram)}]{\sk{अन्तरम्} (antaram) — a gap; the distance between}

\label{lesson:SA-C423-antaram}



\begin{quote}
Before the new one: say the Sanskrit for an exit, then the Sanskrit for opposite.
\end{quote}

\subsection*{You'll want to know: \sk{अन्तरम्}}
\textbf{\sk{अन्तरम्}} — \emph{antaram} — \textquotedblleft{}a gap; the distance between\textquotedblright{}.

\begin{quote}
\sk{गृहस्य} \sk{मार्गस्य} \sk{च} \sk{अन्तरम्} \sk{अल्पम्} \sk{अस्ति।}
\end{quote}

\begin{quote}
\emph{gṛhasya mārgasya ca antaram alpam asti.}
\end{quote}

\begin{quote}
The gap between the house and the road is small.
\end{quote}

A neuter noun.

\begin{cousinweb}
\textbf{\sk{अन्तरम्}} (\emph{antaram}) is built on \emph{antar}, \textquotedblleft{}within, between\textquotedblright{}; MW gives \emph{antara} as interior and, as a noun, an interval, a distance, a gap (MW).
\end{cousinweb}

\begin{grammarlens}[title={What you've built}]
One more word for getting around.
\end{grammarlens}

\subsection*{Your turn}
\begin{itemize}
\raggedright
  \item \emph{Say it:} \emph{antaram}
  \item \emph{Say it:} \emph{antaram}, once more
  \item \emph{Say it:} \emph{gṛhasya mārgasya ca antaram alpam asti.} — The gap between the house and the road is small.
  \item \emph{Recall:} say the Sanskrit for an exit, then the Sanskrit for opposite, then say \emph{antaram} again
\end{itemize}

\begin{tcolorbox}[breakable,colback=teal!4,colframe=teal!35!black,title={Before you move on}]
What does \sk{अन्तरम्} mean? (\textquotedblleft{}A gap; the distance between\textquotedblright{}.) Say it once more.
\end{tcolorbox}
\section[\texorpdfstring{\sk{द्वारपालः} (dvārapālaḥ)}{द्वारपालः (dvārapālaḥ)}]{\sk{द्वारपालः} (dvārapālaḥ) — a gatekeeper}

\label{lesson:SA-C423-dvarapalah}



\begin{quote}
Before the new one: say the Sanskrit for opposite, then the Sanskrit for a gap.
\end{quote}

\subsection*{You'll want to know: \sk{द्वारपालः}}
\textbf{\sk{द्वारपालः}} — \emph{dvārapālaḥ} — \textquotedblleft{}a gatekeeper\textquotedblright{}.

\begin{quote}
\sk{द्वारपालः} \sk{प्रवेशे} \sk{तिष्ठति।}
\end{quote}

\begin{quote}
\emph{dvārapālaḥ praveśe tiṣṭhati.}
\end{quote}

\begin{quote}
The gatekeeper stands at the entrance.
\end{quote}

A masculine noun.

\begin{cousinweb}
\textbf{\sk{द्वारपालः}} (\emph{dvārapālaḥ}) is built from \emph{dvāra}, \textquotedblleft{}a door, a gate\textquotedblright{}, and \emph{pāla}, \textquotedblleft{}a guard\textquotedblright{}, from the root √pā, \textquotedblleft{}to protect\textquotedblright{}; MW gives the compound as a doorkeeper (MW).
\end{cousinweb}

\begin{grammarlens}[title={What you've built}]
One more word for getting around.
\end{grammarlens}

\subsection*{Your turn}
\begin{itemize}
\raggedright
  \item \emph{Say it:} \emph{dvārapālaḥ}
  \item \emph{Say it:} \emph{dvārapālaḥ}, once more
  \item \emph{Say it:} \emph{dvārapālaḥ praveśe tiṣṭhati.} — The gatekeeper stands at the entrance.
  \item \emph{Recall:} say the Sanskrit for opposite, then the Sanskrit for a gap, then say \emph{dvārapālaḥ} again
\end{itemize}

\begin{tcolorbox}[breakable,colback=teal!4,colframe=teal!35!black,title={Before you move on}]
What does \sk{द्वारपालः} mean? (\textquotedblleft{}A gatekeeper\textquotedblright{}.) Say it once more.
\end{tcolorbox}
